\section{Conclusiones y trabajo futuro}
\label{sec:conclusiones}

\begin{table}[!htb]
\centering\footnotesize
\caption{Limitaciones del trabajo.}
\label{tab:riesgos}
\rowcolors{2}{SoftGray}{white}
\begin{tabularx}{\textwidth}{@{}>{\bfseries\raggedright\arraybackslash}p{5.0cm} >{\raggedright\arraybackslash}X@{}}
\toprule
\rowcolor{TableHead} Limitación & Consecuencia y manejo \\
\midrule
Sensor térmico y tonos simulados en la FPGA & Los bloques que los procesan son los definitivos; falta medir alcance y falsas alarmas \\
No hubo pruebas con insectos & La respuesta al estímulo y la marcha vienen de la literatura \\
El radio es un cable & El alcance de 32~m es un parámetro; la atenuación está calculada, no medida \\
Posición conocida de cada mochila & En el sistema real saldría de la navegación inercial y de los tonos \\
El receptor solo escucha & Si A y B quedan a 32~m o más entre sí, solo el receptor muestra el total \\
Versión final de la red sin repetir en tarjetas & Verificada por simulación y síntesis; en hardware se probó la versión anterior \\
Colapsos sintéticos y en dos dimensiones & Sirven para comparar políticas, no para predecir tiempos de rescate \\
Ambiente caliente y confinamiento legal & Límites del animal y de la ley: definen dónde no debe usarse \\
\bottomrule
\end{tabularx}
\end{table}

\begin{enumerate}
  \item \textbf{El insecto carga menos de lo que se cita.} La meta defendible es una mochila de 2 a 3~g y de perfil
        bajo; la masa del cómputo es despreciable y mandan la batería y el sensor.
  \item \textbf{El núcleo existe, está verificado y cabe en un chip.} 477 elementos lógicos, idéntico a su modelo de
        referencia en 11\,424 cuadros y probado en dos tarjetas. En SKY130 ocupa 0.050~mm$^2$, pasa DRC y LVS y
        alcanza la misma frecuencia que en la FPGA, con 0.58~mA a 1.8~V en espera.
  \item \textbf{Tres FPGA bastan para detectar, avisar y ubicar sin computadora.} La red retransmite el aviso de la
        mochila fuera de alcance y triangula con 6.2~cm de error a 31.95~m, en 1.10~s. Pasa los nueve escenarios,
        incluidos el ruido y el corte exacto en 32~m.
  \item \textbf{En punto fijo, el error se mide, no se supone.} Dieciséis iteraciones del CORDIC dan 0.017\textdegree{} y
        más no mejoran; la distancia por CORDIC truncado fallaba en más de la mitad de las distancias exactas y una
        raíz entera la dejó sin error con la quinta parte de la lógica; el PID difiere del de aritmética real en menos
        de una cuenta.
  \item \textbf{Guiar ayuda, sin milagros.} La detección en una hora sube de \pPasivaSesenta~\% a \pGuiadaSesenta~\%
        con diez insectos, pero pesan más el número de insectos y la forma del colapso, y en uno de los ocho guiar
        empeora el resultado.
  \item \textbf{Para que vuelvan hace falta una señal en la entrada.} El refugio recupera \regTrampaDos~\% en dos horas
        y un faro, \regFaroDos~\%; ningún método los recupera a todos.
\end{enumerate}

\textbf{Sobre la hipótesis.} Se sostiene en parte: la política aumenta la detección sin habituación y cabe en
miligramos de silicio, pero la mejora depende del colapso. Que dos mochilas bastan para ubicar a la persona se
cumplió en simulación y síntesis con el sensor simulado; falta repetirlo con sensores reales.

\textbf{Trabajo futuro.} (1) Enviar el GDS a una corrida compartida como Tiny Tapeout~\cite{tinytapeout}, medirlo en
silicio y llevar al mismo flujo el detector de tonos. (2) Conectar una matriz AMG8833 o MLX90640, un micrófono MEMS y
un sensor de CO$_2$. (3) Estimar la distancia por tonos, para que la posición deje de darse por conocida. (4) Más
nodos y un receptor que retransmita. (5) El regreso con faro y dos fotodiodos, cerrando el lazo con el PID ya
verificado. (6) Pruebas en recinto cerrado con la colonia de la universidad y el protocolo de la
sección~\ref{sec:etica}.
